\documentclass{article}
\usepackage[T1]{fontenc}
\usepackage{lmodern}
\usepackage{graphicx}
\usepackage{amsmath,amssymb}
\usepackage[a4paper,margin=2cm]{geometry}
\usepackage[absolute,overlay]{textpos}
\setlength{\TPHorizModule}{1mm}
\setlength{\TPVertModule}{1mm}
\usepackage[indonesian]{babel}
\usepackage{hyperref}
\setlength{\emergencystretch}{2em}
% unit: Halaman 40

\begin{document}

% === Halaman 40 (llm) ===

\section*{1. RC4}

\begin{itemize}
    \item \textit{Cipher}alir yang paling popular, namun sekarang sudah tidak aman
    \item Dibuat oleh Ronald Rivest (1987) dari Laboratorium RSA
    \item \textit{RC} adalah singkatan dari \textit{Ron's Code}. Versi lain mengatakan \textit{Rivest Cipher}.
    \item Digunakan di dalam beberapa sistem keamanan seperti:
    \begin{itemize}
        \item protokol SSL (\textit{Secure Socket Layer}) dan TLS (\textit{Transport Layer Security})
        \item Standard IEEE 802.11 \textit{wireless} LAN: WEP (\textit{Wired Equivalent Privacy})
        \item WPA (\textit{Wi-fi Protocol Access}) untuk nirkabel
    \end{itemize}
\end{itemize}

\end{document}
